% Pillar-1 (scaling) standalone abstract
\begin{abstract}
Group Relative Policy Optimization (GRPO) has become a default algorithm for
reinforcement-learning (RL) post-training of language models, yet how its benefit
scales with model size remains unclear and is easily confounded by implementation
and evaluation choices. We study GRPO scaling as one pillar of \ours{}, a
benchmark spanning 70+ runs across seven RL libraries and five model families
($0.6$B--$\sim$671B parameters) on GSM8K, HumanEval, and synthetic tool-use tasks.
Crucially, these runs use \emph{managed} and open back-ends rather than local
frontier clusters: the frontier-scale anchors (Qwen3-235B, DeepSeek-V3.1
$\sim$671B, Nemotron-120B) are obtained through the \textsc{Tinker} managed
training/inference API in a single-seed, descriptive-only regime, which is what
makes a cross-scale study tractable without owning frontier hardware (and why full
base weights are not released for those anchors).
Our central result is a conservative negative one: across roughly $2.4$ orders of
magnitude in parameter count spanned by the fitted anchor pool (a subset of the
full $0.6$B--$\sim$671B roster), the cross-scale slope of the mean GRPO training
reward shows no reliable positive trend---flat within the (wide) uncertainty of
these mostly single-seed anchors, which we therefore read \emph{descriptively}
rather than as a powered significance test---and no single saturation exponent is identifiable
(four of the five frontier-scale traces have their saturation rate pinned at the
fit boundary, and none is identifiable). A pre-registered test geometrically falsifies a clean three-phase
saturation hypothesis (only $2$ of $12$ anchors match), and the Nemotron-120B run
is better described as a distinct collapse phase (zero-reward step fraction $0.55$
versus $\le 0.067$ elsewhere) than as a point on a smooth trend. The defensible
object is therefore a local, stack-conditioned taxonomy of endpoint ceilings,
with categorical capability structure carrying more signal than $\log_{10} N$,
not a Chinchilla-style power law in $N$.
\paragraph{Venue claim boundary (workshop-short / limits audit).}
This manuscript is submitted as a \emph{descriptive limits and identifiability
audit}, not as a positive scaling law or a multi-seed cross-scale significance
study. Frontier anchors remain single-seed and descriptive by design; matched
multi-seed cross-scale runs are an explicitly stated \emph{non-claim} for this
submission and do not gate the negative/identifiability result we headline.
We
release all code, logs, per-step telemetry, and trained adapters; full base-model
checkpoints are shared where model licensing permits, while the managed
frontier-scale runs expose logs and telemetry rather than base weights.
\end{abstract}
